\par\medskip
\noindent
11.3. Utilizaremos la expresión explícita (11.2.1) para esbozar la
definición de un acoplamiento canónico
\begin{equation}\tag{11.3.1}
\Phi_o(\ell)\otimes\Phi'_o(\ell)\longrightarrow\mathbb Q/\mathbb Z
\end{equation}
que es una dualidad perfecta; aquí $\Phi'_o$ denota, como en la sección
1, el grupo de componentes conexas asociado al modelo de Néron del
\underline{dual} $A'_K$ de $A_K$. Apenas cabe duda (aunque el redactor
ha prescindido de la verificación) de que este acoplamiento no es otro,
quizá salvo el signo, que el definido en (1.2.1), y se invita al lector,
como ejercicio, a establecer la compatibilidad requerida (¡con el
signo correcto!).

Pongamos
\begin{equation*}
U'=T_\ell\bigl(A'(\bar K)\bigr),
 \qquad
T=T_\ell(1)(\bar K)=T_\ell(\bar K^*).
\end{equation*}
Existe un \underline{acoplamiento perfecto} de
$\mathbb Z_\ell$-módulos libres de tipo finito, compatible con las
acciones de $I$,
\begin{equation}\tag{11.3.2}
\varphi:U\otimes U'\longrightarrow T,
 \qquad (T\simeq\mathbb Z_\ell\text{ no canónicamente}),
\end{equation}

%%%%%%%%%%  scan idx 463  |  folio 452  %%%%%%%%%%

y, por la estructura conocida de $I$, una sucesión exacta
\begin{equation}\tag{11.3.3}
1\longrightarrow R\longrightarrow I\longrightarrow T\longrightarrow1,
\end{equation}
donde $R$ es un grupo profinito \underline{primo con} $\ell$. Partiendo
ahora de los datos (11.3.2) y (11.3.3), introduzcamos los grupos
\begin{equation}\tag{11.3.4}
\begin{aligned}
\Phi&=\frac{(U\otimes D_\ell)^I}{U^I\otimes D_\ell},\qquad
\Phi'=\frac{(U'\otimes D_\ell)^I}{{U'}^I\otimes D_\ell},\qquad
D_\ell=\mathbb Q_\ell/\mathbb Z_\ell.
\end{aligned}
\end{equation}
Definiremos una dualidad perfecta
\begin{equation}\tag{11.3.5}
\Phi\otimes\Phi'\longrightarrow\mathbb Q/\mathbb Z,
\end{equation}
que dará, en particular, el acoplamiento anunciado (11.3.1).
Consideremos la sucesión exacta de $I$-módulos
\begin{equation}\tag{11.3.6}
0\longrightarrow U\xrightarrow{\ i\ }
 U\otimes_{\mathbb Z_\ell}\mathbb Q_\ell
 \xrightarrow{\ j\ }U\otimes D_\ell\longrightarrow0
\end{equation}
y la sucesión exacta de cohomología asociada, donde la cohomología de
$I$ se define mediante \underline{cocadenas continuas} con coeficientes
en los grupos topológicos considerados:
\begin{equation}\tag{11.3.7}
\resizebox{\textwidth}{!}{%
\begin{tikzcd}[ampersand replacement=\&, column sep=small, row sep=large]
0 \arrow[r]
 \& H^o(I,U) \arrow[r, "i^o"] \arrow[d, no head, "\wr"]
 \& H^o(I,U\otimes\mathbb Q_\ell) \arrow[r, "j^o"]
      \arrow[d, no head, "\wr"]
 \& H^o(I,U\otimes D_\ell) \arrow[r, "\partial"]
      \arrow[d, no head, "\wr"]
 \& H^1(I,U) \arrow[r, "i^1"]
 \& H^1(I,U\otimes\mathbb Q_\ell) \arrow[d, no head, "\wr"] \\
{}
 \& U^I
 \& (U\otimes\sisi{Q_\ell}{\mathbb Q_\ell})^I\simeq U^I\otimes_{\mathbb Z}\mathbb Q
 \& (U\otimes D_\ell)^I
 \& {}
 \& H^1(I,U)\otimes_{\mathbb Z}\mathbb Q .
\end{tikzcd}}
\end{equation}
Es inmediato que el grupo $\Phi$ definido por (11.3.4) es precisamente
$\operatorname{coker}j^o$ y, por tanto, es canónicamente isomorfo a
$\ker i^1$. Obtenemos así una sucesión exacta canónica
\begin{equation}\tag{11.3.8}
0\longrightarrow\Phi\longrightarrow H^1(I,U)
 \longrightarrow H^1(I,U)\otimes_{\mathbb Z}\mathbb Q,
\end{equation}

%%%%%%%%%%  scan idx 464  |  folio 453  %%%%%%%%%%

que muestra que $\Phi$ \underline{es canónicamente isomorfo al
subgrupo de torsión de $H^1(I,U)$}.

Por otra parte, de la sucesión exacta (11.3.3) se deduce, para todo
$n=\ell^r$, un isomorfismo
\begin{equation}\tag{11.3.9}
H^1(I,T_n)\simeq\operatorname{Hom}(I,T_n)
 \xleftarrow{\sim}\operatorname{Hom}(T,T_n)
 \xleftarrow{\sim}\mathbb Z/n\mathbb Z.
\end{equation}
(Obsérvese que $I$ actúa trivialmente sobre $T$ y, por tanto, sobre
$T_n=T/nT$.) Se sigue que, para todo $I$-$\mathbb Z_\ell$-módulo finito
$M$, si
\begin{equation}\tag{11.3.10}
M'=\operatorname{Hom}(M,T\otimes D_\ell)
\end{equation}
denota su ``dual de Cartier'', existe un acoplamiento canónico
\begin{equation}\tag{11.3.11}
H^i(I,M)\otimes H^{1-i}(I,M')
 \longrightarrow H^1(I,T\otimes D_\ell)\simeq D_\ell,
\end{equation}
donde el último isomorfismo se deduce de los isomorfismos (11.3.9)
pasando al límite inductivo sobre las potencias $n$ de $\ell$. Es bien
sabido que estos acoplamientos son \underline{dualidades perfectas}
(véase, por ejemplo, SGA 5 I 5). Tomando en estos acoplamientos $M$ de
la forma $U_n$, de modo que $M'=U'_n$, y pasando al límite, obtenemos un
acoplamiento perfecto
\begin{equation}\tag{11.3.12}
H^o(I,U\otimes D_\ell)\times H^1(I,U')\longrightarrow D_\ell,
\end{equation}
donde el segundo factor del miembro izquierdo es un
$\mathbb Z_\ell$-módulo de tipo finito. Un nuevo paso al límite en los
acoplamientos (11.3.12), para morfismos de transición de la forma
$\ell^r\operatorname{id}$, da un segundo acoplamiento perfecto
\begin{equation}\tag{11.3.13}
H^o(I,U\otimes\mathbb Q_\ell)
 \times H^1(I,U'\otimes\mathbb Q_\ell)
 \longrightarrow D_\ell
\end{equation}
de espacios vectoriales de dimensión finita sobre $\mathbb Q_\ell$.
Denotemos por $i'$ y $j'$ los

%%%%%%%%%%  scan idx 465  |  folio 454  %%%%%%%%%%

morfismos definidos como en (11.3.6), con $U$ sustituido por $U'$. Es
entonces fácil ver que los homomorfismos siguientes
\begin{center}
\begin{tikzcd}[column sep=6em, row sep=5em]
H^o(I,U\otimes\mathbb Q_\ell) \arrow[r, "j^o"]
  & H^o(I,U\otimes D_\ell) \\
H^1(I,U'\otimes\mathbb Q_\ell)
  & H^1(I,U') \arrow[l, "{i'}^1"']
\end{tikzcd}
\end{center}
son \underline{traspuestos} uno de otro respecto de los acoplamientos
(11.3.12) y (11.3.13). Se deduce inmediatamente que
$\operatorname{coker}j^o$ y $\ker {i'}^1$ son duales entre sí mediante
un acoplamiento canónico con valores en $D_\ell$; este es el
acoplamiento anunciado (11.3.5).
